\section{内核和用户程序怎样约定二进制接口}

内核和用户程序分别编译，却要对系统调用号、寄存器、结构大小和字段 offset
得出相同结论。项目把这些共享定义集中到 ABI 头文件，并用静态断言、布局测试
和真实用户 ELF 检查漂移。devfs 与 procfs 也在这一阶段接入 VFS：前者提供
设备端点，后者在读取时生成内核状态。

\topic{增加一个系统调用会同时改动哪些位置}
假设在调用号 30 和 31 之间插入一个新调用。重新编译后的 Kernel 会把旧用户
程序发来的 31 解释成另一项操作；编译仍然成功，运行结果却已经改变。若参数
还是一个结构体，字段顺序、宽度、对齐和错误值也必须由内核与旧 ELF 作出相同
解释。

因此这一章不再插入调用号，而是做三件具体的事：

\begin{enumerate}[leftmargin=2.1em]
  \item 把已经存在的用户/内核约定写成共享且可编译的 ABI；
  \item 用现有 VFS 契约补齐设备命名和只读观察面；
  \item 用已有系统调用组合用户工具，不为“工具数量”发明新的内核机制。
\end{enumerate}

整理这些边界仍然会发现代码错误。例如 \code{cp} 若先截断目标、再发现源和
目标是同一 inode，源数据已经丢失；procfs 若持自己的锁调用 VFS 快照，两把
锁可能形成环；多个 TCG 实例并行运行，也可能让正确用例撞上墙钟超时。

\topic{ABI 的历史：源代码兼容与二进制兼容不是一回事}
\term{API}{Application Programming Interface}描述源代码怎样调用一个
接口。头文件变了，调用者与实现一起重新编译，很多变化仍可被吸收。
\term{ABI}{Application Binary Interface}约束的是已经生成的机器码怎样
交换信息。磁盘上的旧 ELF 不会因为 Kernel 头文件改变而自动重编译，因此
下列事实都会成为协议：

\begin{itemize}[leftmargin=2em]
  \item 系统调用号，以及参数位于哪些 64 位寄存器；
  \item 成功返回和负错误值的解释；
  \item 结构大小、对齐、字段 offset 与枚举底层值；
  \item 指针宽度、字节序和用户地址验证规则；
  \item ELF class、machine、文件类型、段权限和入口条件；
  \item signal frame 这类由 Kernel 写入、再由用户态交回 Kernel 的双向对象。
\end{itemize}

Unix 家族长期兼容的历史说明了 ABI 的代价：旧 syscall number、结构变体和
设备控制命令一旦进入程序，就很难仅靠“设计得更整洁”而删除。本项目不复制
Linux ABI，但学习同一条工程规律。ABI major 为 2 表示这一代边界；minor
和 patch 只允许在旧消费者仍可正确解释的条件下演进。破坏字段或语义必须
提高 major，不能静默复用旧号。

\begin{contractbox}
ABI v2.0.0 固定 69 个系统调用，合法调用号从 1 到 69，公共错误区间从
\(-1\) 到 \(-57\)。目标程序必须是小端 ELF64、x86-64 的
\code{ET_EXEC}。syscall number、errno、结构布局和调用约定均由本项目定义，
与 Linux ABI 不兼容。
\end{contractbox}

\topic{固定宽度类型只解决位数，不自动解决序列化}
宿主环境中的 \code{long} 在 LP64 和 LLP64 下宽度不同，\code{size_t}
表达当前编译目标可容纳的对象大小，\code{bool} 的外部表示也不适合成为线
协议。ABI、磁盘布局、硬件寄存器和汇编交接因此只使用
\code{uint8_t}、\code{uint16_t}、\code{uint32_t}、
\code{uint64_t} 及对应有符号类型。

然而固定宽度并不等于完成协议设计。它没有自动决定大小端，没有消除结构
padding，也不会验证枚举是否合法、地址相加是否溢出、指针是否映射或页面
是否可写。于是三类边界仍采用不同策略：

\begin{itemize}[leftmargin=2em]
  \item rootfs 和 journal 逐字段执行显式小端 Store/Load，不把 C++ 对象
        原样写盘；
  \item syscall 结构先从用户页 copy-in，再验证全部字段，成功后才提交状态；
  \item ELF 先验证 header 与 program header 的范围、权限和重叠，再读取段。
\end{itemize}

“用大类型为未来留余量”只适用于容量和计数语义本身确实需要 64 位的场景，
不能用来隐藏协议不清晰。地址、文件长度、inode、时间和累计统计选择
\code{uint64_t}；磁盘扇区中的固定位域仍按盘面定义选择准确宽度。

\topic{编译器怎样成为 ABI 检查器}
\filepath{source/abi/include/os/abi/layout.hpp} 对共享结构执行
\code{offsetof} 与 \code{static_assert}。检查总大小而不检查 offset
是不够的：两个结构可能总大小相同，字段顺序却已经互换。只检查 offset
也不够：枚举号、错误边界或目标 ELF 身份仍可能漂移。

\begin{osgridlongtable}{L{0.28\linewidth}L{0.22\linewidth}L{0.38\linewidth}}
对象 & 机器码依赖的值 & 改动后会发生什么 \\ \hline
\endfirsthead
对象 & 机器码依赖的值 & 改动后会发生什么 \\ \hline
\endhead
\code{ProcessLaunchRequest} &
六个 64 位字段位于 0、8、16、24、32、40 &
Kernel 必须用与用户程序相同的方式解释路径、参数和环境 \\ \hline
\code{SignalFrame} &
header 0--56，\code{context} 从 64 开始 &
signal handler 返回时会把同一对象再次交给 Kernel 验证 \\ \hline
\code{FileInformation} &
身份、类型、逻辑/分配长度和链接数的 offset &
\code{stat}、\code{cp} 与工具不能把长度解释成 inode \\ \hline
\code{TerminalInformation} &
terminal、session、foreground group 各 64 位 &
TTY 作业控制的权限判断横跨 Ring 0 与 Ring 3 \\ \hline
\code{SystemCallNumber} &
首项、末项与总数 &
中间插值或重排不能在编译仍成功时改变旧程序语义 \\ \hline
ELF 常量 &
magic、class、data、machine、flags &
Stage 1、Kernel 和产物审计必须接受同一种目标身份 \\ \hline
\end{osgridlongtable}

共享断言由 Kernel 与用户程序共同编译，独立
\filepath{tests/unit/abi_v2_contract_test.cpp} 再从消费者角度检查版本、
枚举首尾、结构大小、对齐和全部跨边界字段。前者让错误尽早停在编译阶段，
后者防止某个目标没有真正包含共享头或测试只覆盖了局部结构。

\begin{failurebox}
不要在 ABI 头中加入 C++ 容器、虚函数对象、宿主指针或依赖实现的位域。
这些对象即使在当前编译器上“看起来能用”，也没有稳定的外部对象表示。ABI
结构只表达数据；所有权、锁和行为留在各自地址空间的实现对象中。
\end{failurebox}

\topic{ELF 常量为何也属于 ABI}
此前 Stage 1、Kernel 用户 ELF loader 与离线审计都需要知道
\code{0x7F ELF}、64 位 class、小端 encoding、x86-64 machine 和段权限。
若三处各有一组魔法数字，它们可能逐渐形成三个略有差异的“标准”。代码
把目标身份集中到 \filepath{source/abi/include/os/abi/elf.hpp}，Kernel 的
\filepath{source/kernel/src/user/user_elf.cpp} 直接消费它。

共享常量并不意味着复用同一份 parser。Stage 1 仍在受限启动环境中解析
Kernel ELF，Kernel 仍对不可信用户 ELF 进行两遍事务加载，Python 审计仍是
独立 oracle。应当共享的是协议事实，而不是把生产 parser 原样调用一遍并
宣布“双方结果一致”。

\section{devfs、procfs 与用户工具}

为了让 fd 0/1/2 指向终端，早期实现曾用专用后端提供
\filepath{/dev/console}。这让 open/read/write 可以走通，却把两个责任粘在一起：

\begin{enumerate}[leftmargin=2.1em]
  \item \code{console} 这个名称怎样出现在 \filepath{/dev}；
  \item 打开后字符怎样由 Terminal/TTY 读取和写出。
\end{enumerate}

真实 Unix 系统把设备表示成文件节点，历史原因在于用户程序可以复用
\code{open/read/write/close}，而不必为每种设备发明一套调用。devfs 学习的
正是“命名空间与驱动分层”，不是模拟 Linux 的 major/minor、udev、热插拔
和权限数据库。

\begin{pdfdiagram}
  \node[os state] (path) {VFS path\\\filepath{/dev/console}};
  \node[os state,right=of path] (devfs) {devfs\\名称与 vnode};
  \node[os block,right=of devfs] (fd) {FileDescription\\对象类型};
  \node[os hardware,below=of fd] (tty) {Terminal / TTY\\字符与作业控制};
  \draw[os arrow] (path) -- node[above]{lookup} (devfs);
  \draw[os arrow] (devfs) -- node[above]{open} (fd);
  \draw[os arrow] (fd) -- node[right]{read/write} (tty);
\end{pdfdiagram}
\diagramnote{devfs 不保存驱动指针；路径、打开对象和字符行为各自只有一个所有者。}

\topic{devfs 的注册、发布与 vnode 生命周期}
\code{Devfs} 由调用者提供固定 \code{DevfsDevice[16]}。每项只保存稳定
node identifier、generation、名称长度、最多 255 字节名称和 active 状态。
初始化后注册 \code{console}，再把 superblock 挂载到 \filepath{/dev}。

注册在锁内完成一次完整扫描。扫描既检查 active 项是否重名，也记住第一个
空槽；只有名称、容量和重名检查全部成功后，才一次性发布完整新项。绝不能先
置 active 再复制名字，否则并发 lookup 可能观察到半个设备节点。

\begin{lstlisting}[language=C++,caption={devfs 注册的事务顺序（概念代码）}]
ValidateName();
Lock();
FindDuplicateAndFirstFreeSlot();
if (duplicate || no_free_slot) {
    RecordRejection();
    Unlock();
    return failure;
}
FillNodeIdGenerationName();
PublishActive();
Unlock();
\end{lstlisting}

当前版本没有注销，因此已经返回的 vnode 不会因为槽位重用而突然指向另一
设备。未来若加入 hot unplug，必须同时设计 generation 增长、打开引用、
正在阻塞的 I/O、设备消失错误和槽位复用顺序，不能只增加一个
\code{Unregister} 函数。

devfs 命名空间只读。create、remove、rename 和 truncate 返回只读错误；
普通文件 read/write 返回不支持。真正的字符 I/O 由 FileDescription 选择
TerminalDevice，不经 devfs 普通文件回调。这样同一设备不会拥有两条内容
不一致的读写路径。

\topic{procfs 的历史动机与本项目范围}
“一切皆文件”从来不是说所有对象都真的位于磁盘，而是说不同资源尽量复用
统一的名称和 I/O 协议。早期 Unix 的 \filepath{/dev} 让设备进入文件 API；
后来 procfs 让进程与系统状态也可由普通工具读取。Linux 的 procfs 经过多年
增长，包含大量历史字段和可写控制项，并非一套小而严格的永久 ABI。

本项目只取其中最适合教学的一层：以 VFS 暴露动态只读观察面。固定六个文件：

\begin{osgridlongtable}{L{0.24\linewidth}L{0.62\linewidth}}
路径 & 一次快照中的含义 \\ \hline
\endfirsthead
路径 & 一次快照中的含义 \\ \hline
\endhead
\filepath{/proc/version} & ABI major/minor 与 \code{x86_64} 目标架构 \\ \hline
\filepath{/proc/uptime} & PIT 驱动的单调纳秒，不是宿主墙钟 \\ \hline
\filepath{/proc/meminfo} & managed、free、allocated 物理内存字节 \\ \hline
\filepath{/proc/processes} & 活动 Process/Thread、容量与当前 PID \\ \hline
\filepath{/proc/resources} & heap、FileDescription、Pipe、vnode 与 journal commit \\ \hline
\filepath{/proc/mounts} & 当前 VFS mount 数 \\ \hline
\end{osgridlongtable}

没有每 PID 目录、\filepath{/proc/self}、符号链接、权限、命令行和可写调优。
当前没有这些机制；加入它们需要新的 vnode 类型、权限字段和生命周期处理。

\topic{为什么每次 read 都要重新生成快照}
进程数、空闲页和单调时间会持续变化。若挂载时生成文本，内容在第一次读取前
就可能过时；若每个 open 长期保存完整快照，则需要动态分配、新的关闭生命
周期和更复杂的偏移状态。因此每次 read 或 stat 都重新采集一次
\code{ProcfsSnapshot}，使用 256 字节局部缓冲格式化。

\begin{pdfdiagram}
  \node[os state] (read) {VFS read\\offset/capacity};
  \node[os block,right=of read] (capture) {采集 15 个\\\code{uint64_t}};
  \node[os block,right=of capture] (render) {有界十进制\\局部缓冲};
  \node[os state,below=of render] (slice) {short read / EOF\\复制请求切片};
  \draw[os arrow] (read) -- (capture);
  \draw[os arrow] (capture) -- (render);
  \draw[os arrow] (render) -- (slice);
\end{pdfdiagram}
\diagramnote{一次 read 内使用同一快照；下一次 read 可以看到新的系统状态。}

offset 大于等于渲染长度时返回 EOF；否则只复制
\(\min(\text{capacity}, \text{size}-\text{offset})\)。stat 也重新渲染以给出
当前文本长度，所以 stat 与随后 read 之间状态变化可能导致长度不同。动态
伪文件的用户程序必须正确处理 short read 和 EOF，不能把 stat size 当成
跨调用事务承诺。

\topic{锁顺序：最容易被文本文件隐藏的内核死锁}
snapshot provider 需要读取物理页、heap、Process/Thread、FileDescription、
Pipe、VFS、rootfs journal 和 PIT 时间。若 procfs 先持自己的锁再调用这些
子系统，而某个资源统计路径又反向访问 procfs，就形成锁环。

\begin{failurebox}
禁止在持有 procfs 统计锁时调用 snapshot provider、VFS 或其他子系统。
procfs 锁只保护自己的 active open、成功读取、失败和累计字节统计。快照
采集与文本格式化期间不持该锁；完成后再短暂加锁更新有界计数。
\end{failurebox}

\begin{pdfdiagram}
  \node[os state] (entry) {procfs read};
  \node[os block,right=of entry] (provider) {无 procfs 锁\\调用 provider};
  \node[os block,right=of provider] (systems) {memory / process\\VFS / journal};
  \node[os state,below=of provider] (stats) {短暂持锁\\更新统计};
  \draw[os arrow] (entry) -- (provider);
  \draw[os arrow] (provider) -- (systems);
  \draw[os arrow] (systems) |- (stats);
\end{pdfdiagram}
\diagramnote{跨子系统读取放在锁外，procfs 自身统计放在末尾，锁方向不会倒置。}

VFS 汇总也不能在全局锁内递归回调后端。当前 mount 数组在初始化完成后不再修改，
汇总逐个调用后端的 \code{ReadResourceUsage}，并在每次加法前检查
\code{uint64_t} 溢出。任何后端失败或总和溢出都清空输出并返回失败，不能
提供“看起来合理”的部分统计。procfs 自己的 resource usage 只报告固定节点，
不再次采集快照，因此不会递归进入自身。

\topic{没有 libc 时怎样安全输出十进制}
freestanding Kernel 不依赖 \code{snprintf}、iostream 或动态 string。
\code{SnapshotWriter} 用最多 20 字节临时区保存一个 \code{uint64_t}：
反复取模 10 得到低位、除以 10 推进、再反转到正序。零必须单独输出一个字符，
最大 64 位无符号数恰好 20 位。

每次追加先验证 \code{length <= capacity - size}，而不是直接判断
\code{size + length <= capacity}；后者的加法本身可能先溢出。任一字段无法
完整写入时，整个 render 返回容量耗尽，不把截断文本冒充成功文件。

累计统计同样采用饱和加法。观察计数到达 \code{uint64_t} 最大值后保持最大值，
而不是回绕为零。active open 则在达到最大值时拒绝新 open，因为继续增加会
使 active open 无法与 close 次数相互核对。

\topic{32 个工具怎样只用现有系统调用完成工作}
rootfs 安装 32 个独立工具路径：

\begin{lstlisting}[caption={rootfs 中安装的用户工具}]
help echo cat wc head tee true false pwd ls stat mkdir write touch
rm rmdir mv truncate sync basename dirname cp seq uptime ps free uname
mounts resources sleep kill id
\end{lstlisting}

它们可以共享 multi-call ELF 代码字节，但每个安装路径拥有独立 inode。
\code{basename} 和 \code{dirname} 只处理用户字符串；
\code{uptime/ps/free/uname/mounts/resources} 只读取 procfs；
\code{cp} 组合 stat、open、read、write 与 close；\code{sleep} 使用已有
deadline；\code{kill} 使用已有 signal；\code{id} 查询当前 PID。

工具同样需要有界失败语义。\code{seq} 最多输出 100000 个数；十进制解析在
加入每个 digit 前检查 \code{value <= (max-digit)/10}；sleep 在毫秒乘以
一百万前检查上限。这样一个小命令不会成为无界串口日志或整数回绕入口。

\begin{failurebox}
\code{cp source source} 不能先以截断方式打开目标。实现先 stat 两端，比较
mount identifier、superblock identifier、inode 与 generation；确认不是
同一文件后才打开目标。路径字符串不同不代表对象不同，\filepath{/a/../b}
和 \filepath{/b} 可能解析到同一 vnode。
\end{failurebox}

\topic{tool probe 为什么不能被普通编译替代}
CMake 能证明目标被列入构建，却不能证明 rootfs 安装后仍有 32 个正确路径。
所有名字若意外指向同一个 inode，或者某个路径没有写入镜像，编译和链接都
可能成功。

\filepath{/bin/tool_probe} 在真实 Ring 3 启动序列中逐项执行：

\begin{enumerate}[leftmargin=2.1em]
  \item stat 成功且类型是 regular file；
  \item 32 个 inode 两两不同；
  \item 文件长度至少容纳 ELF magic；
  \item open/read 得到 \code{7F 45 4C 46}；
  \item 所有 fd 在成功和失败路径都关闭。
\end{enumerate}

全部通过才打印 \code{ELF32_VERIFIED}。256 MiB functional QEMU 随后从
Shell 实际运行新增工具：basename/dirname/seq 使用唯一输出，procfs 工具
检查字段，cp 由 cat 回读，sleep 走真实 deadline，kill 向 PID1 发送默认
忽略的 SIGCHLD。前者证明镜像产物，后者证明 dispatch、参数、syscall、
procfs 和资源回收，二者不能互相替代。

\topic{每类测试拦截哪一种错误}\par
\begin{osgridlongtable}{L{0.23\linewidth}L{0.62\linewidth}}
检查位置 & 拒绝的错误 \\ \hline
\endfirsthead
检查位置 & 拒绝的错误 \\ \hline
\endhead
ABI 单元测试 & 版本、号值、枚举、大小、对齐、字段 offset 或 ELF 身份漂移 \\ \hline
devfs 单元测试 & 重名、容量、半发布、readdir、只读语义、open/close 账本错误 \\ \hline
procfs 单元测试 & 六文件格式、short read、EOF、stat、provider 失败与统计错误 \\ \hline
VFS 集成测试 & rootfs、memfs、devfs、procfs 同时挂载后的路径穿越错误 \\ \hline
随机测试 & 十万组变化数值、offset 与 capacity 相对独立 oracle 的偏差 \\ \hline
产物审计 & Kernel/用户 ELF 段、权限、符号或 freestanding 依赖异常 \\ \hline
64 MiB QEMU & 启动、ABI/devfs/procfs 标记和 32 个独立 ELF 路径 \\ \hline
256 MiB QEMU & 新工具、16 级流水线、TTY、signal 与资源守恒 \\ \hline
64 GiB QEMU & 主规格地址宽度、4 GiB 以上真实页和容量模型 \\ \hline
\end{osgridlongtable}

宿主测试可以并行运行，\code{os_qemu_*} 共享同一个 CTest
\code{RESOURCE_LOCK}。这条锁处理机器资源约束：多个 TCG 实例争抢 CPU 和
内存，会让正确的 2--20 秒用例撞上有界超时。把串行要求写进测试图后，执行者即使使用
\code{ctest -j20} 也不会制造虚假失败。

\topic{日志中的时间与来宾中的时间}
Kernel 只新增三个低频阶段标记：

\begin{lstlisting}
[OS][KERNEL] DEVFS_READY=/dev
[OS][KERNEL] PROCFS_READY=/proc
[OS][KERNEL] ABI_V2_FROZEN
\end{lstlisting}

热路径不逐 open、逐 read 或逐工具打印内核日志，而是在结束时输出有界统计。
QEMU runner 给每行串口加的 \code{[QEMU][T+...ms]} 是宿主收到字节的相对
墙钟，用于判断停滞发生在哪一阶段；\filepath{/proc/uptime} 与内核 deadline
来自来宾 PIT 单调纳秒。前者是测试观测，后者是 OS 时间语义，不能相互冒充。

\topic{源码走读：从协议到真实整机}
建议按“定义、实现、消费者、反例证据”的顺序阅读：

\begin{enumerate}[leftmargin=2.1em]
  \item \filepath{source/abi/include/os/abi/version.hpp}、
        \filepath{elf.hpp} 与 \filepath{layout.hpp}：先看机器码共享哪些值；
  \item \filepath{source/kernel/include/os/kernel/fs/devfs.hpp} 与
        \filepath{source/kernel/src/fs/devfs.cpp}：看注册表怎样发布 vnode；
  \item \filepath{source/kernel/include/os/kernel/fs/procfs.hpp} 与
        \filepath{source/kernel/src/fs/procfs.cpp}：看 snapshot、渲染和锁范围；
  \item \filepath{source/kernel/src/core/kernel_main.cpp}：看 rootfs、memfs、
        devfs、procfs 的挂载顺序和 snapshot provider 怎样接入资源统计；
  \item \filepath{source/user/programs/core_tool.cpp} 与
        \filepath{tool_probe.cpp}：看同一 ABI 的两类真实消费者；
  \item \filepath{tests/unit/abi_v2_contract_test.cpp}、
        \filepath{tests/integration/boot/pseudo_file_system_test.cpp} 与
        \filepath{tests/randomized/kernel/procfs_randomized_test.cpp}：最后看
        独立 oracle 怎样主动构造反例。
\end{enumerate}

走读时对每个函数持续追问：输入来自哪个信任域；加法、乘法和 offset 在哪里
验证；失败前是否发布了半状态；持有哪些锁；对象由谁关闭；测试能否在生产
路径上观察到这些事实。代码命名和注释应使答案能直接从相邻定义获得，而不是
依靠作者记忆。

\topic{ABI、devfs 与 procfs 还缺少什么}
devfs 没有权限、uid/gid、major/minor、注销、热插拔或 udev；procfs 没有
每进程目录、符号链接、可写控制项和 Linux 文本兼容；ABI v2 没有动态链接、
共享库、能力协商或向后加载 v1 用户程序；调度仍是单 BSP、内核不可抢占，
这一实现只适用于当前单核配置；SMP 需要补充跨 CPU 同步。

若以后引入 SMP、权限、动态链接或设备热插拔，旧 ELF 仍会按 ABI v2 的调用号
和结构布局访问内核。兼容实现要保留旧解释；不兼容实现则必须提高 ABI major，
让装载器在运行程序以前明确拒绝。新增对象还要补充销毁顺序、锁顺序和失败后
保留的状态。
